\documentclass{article}
\usepackage{amsmath, amssymb, amsthm}
\usepackage{hyperref}
\usepackage{geometry}
\geometry{margin=1in}

\title{Erd\H{o}s Problem \#970}
\date{Last updated: 2025-08-31}
\author{Source: erdosproblems.com}

\begin{document}
\maketitle

\noindent\textbf{Status:} OPEN \quad \textbf{No prize}

\noindent\textbf{Tags:} number theory

\medskip
\noindent\textbf{Note:} Jacobsthal's function

\medskip

\noindent\textbf{Problem Statement:}

Let $h(k)$ be Jacobsthal's function, defined to as the minimal $m$ such that, if $n$ has at most $k$ prime factors, then in any set of $m$ consecutive integers there exists an integer coprime to $n$. Determine the order of magnitude of $h(k)$. In particular, is it true that\[h(k) \ll k^2?\]




That $h(k)\ll k^2$ is a conjecture of Jacobsthal. Iwaniec \cite{Iw78} proved\[h(k) \ll (k\log k)^2.\]The best lower bound known is\[h(k) \gg \frac{(\log k)(\log\log\log k)}{(\log\log k)^2}k,\]due to Ford, Green, Konyagin, Maynard, and Tao \cite{FGKMT18}. This is a more general form of the function considered in [687] .




References

[FGKMT18] Ford, Kevin and Green, Ben and Konyagin, Sergei and Maynard, James and Tao, Terence, Long gaps between primes . J. Amer. Math. Soc. (2018), 65-105.

[Iw78] Iwaniec, Henryk, On the problem of {J}acobsthal . Demonstratio Math. (1978), 225--231.

\noindent\textbf{Related OEIS sequences:} A048669


\bigskip
\noindent\small{Source: \url{https://www.erdosproblems.com/970}}
\end{document}
